\section{Related Work}

Reinforcement Learning from Human Feedback (RLHF) \cite{10.5555/3294996.3295184}\cite{Ji2024ReinforcementLF} traditionally relies on Proximal Policy Optimization (PPO) \cite{schulman2017proximalpolicyoptimizationalgorithms}, which suffers from heavy computational overhead due to separate reward models. Group Relative Policy Optimization (GRPO) \cite{shao2024deepseekmathpushinglimitsmathematical} circumvents this bottleneck by substituting the critic with a baseline derived from the average rewards of $G$ distinct rollouts. This memory-efficient paradigm was further refined by the Advantage Calibration mechanism introduced in \cite{nan2025ngrponegativeenhancedgrouprelative}, which incorporates virtual maximum rewards to maintain continuous gradient flow in extreme edge cases. Concurrently, augmenting LLM reasoning via Chain-of-Thought \cite{10.5555/3600270.3602070} and autonomous self-correction \cite{madaan2023selfrefineiterativerefinementselffeedback} has gained significant traction. Building on this, Multi-Layer GRPO (MGRPO) \cite{ding2025multilayergrpoenhancingreasoning} employs a two-layer pipeline that utilizes initial outputs as context to verify and correct errors. However, such unconstrained self-reflective frameworks often induce a pronounced correction bias, coercing models to arbitrarily alter valid paths and thereby degrading inference stability.

To mitigate reasoning instability and overcorrection, recent studies focus on uncertainty quantification and sophisticated data curation. Semantic Entropy (SE) \cite{kuhn2023semanticuncertaintylinguisticinvariances} provides an information-theoretic metric to distinguish genuine model confusion from generative diversity-a concept SEED-GRPO \cite{chen2025seedgrposemanticentropyenhanced} leverages to dynamically dampen gradient updates for highly uncertain prompts. Additionally, theoretical insights from Balanced Sampling \cite{gao2026divagrpoenhancingmultimodalreasoning} emphasize that maintaining an equitable distribution of correct and hard-negative samples is pivotal for stabilizing convergence in sparse-reward environments. Collectively, these advancements underscore the necessity of a unified framework. Building upon this foundation, our proposed approach combines semantic error profiling with statically balanced batches to robustly optimize iterative reasoning while circumventing correction bias.
